\acte{Chapelet de saint Michel Archange}
\usage{mar., jour des saints Anges ; 8 mai, 29 sept. ; avant la communion}
\rubrique{Le chapelet de saint Michel Archange, ou chapelet des Anges, se prie sur un chapelet spécial de trente-neuf grains ; les neuf groupes de grains représentent les neuf chœurs des Anges.}
\rubrique{Sur la médaille et les quatre premiers grains :}
\ouverture{S}{aint}{Michel Archange,} grand prince des milices célestes, défendez-nous dans le combat pour que nous ne périssions pas au jour du Jugement.
\vr{Ô Dieu, venez à mon aide.}{Seigneur, hâtez-vous de me secourir.}
\rubrique{Puis : un Pater en l’honneur de saint Michel ; un Pater en l’honneur de saint Gabriel ; un Pater en l’honneur de saint Raphaël ; un Pater en l’honneur de notre Ange gardien.}
\rubrique{Sur chacun des neuf groupes, la salutation, puis un Pater et trois Ave.}
\salut{1}{Par l’intercession de saint Michel et du chœur céleste des Séraphins, que le Seigneur nous rende dignes de brûler du feu de la charité parfaite. Amen.}
\salut{2}{Par l’intercession de saint Michel et du chœur céleste des Chérubins, que le Seigneur nous fasse la grâce d’abandonner la voie du péché et de courir dans celle de la perfection chrétienne. Amen.}
\salut{3}{Par l’intercession de saint Michel et du chœur céleste des Trônes, que le Seigneur insuffle dans nos cœurs un véritable et sincère esprit d’humilité. Amen.}
\salut{4}{Par l’intercession de saint Michel et du chœur céleste des Dominations, que le Seigneur nous donne la grâce de dominer nos sens et de soumettre nos passions indisciplinées. Amen.}
\salut{5}{Par l’intercession de saint Michel et du chœur céleste des Puissances, que le Seigneur daigne protéger nos âmes contre les embûches et les tentations du démon. Amen.}
\salut{6}{Par l’intercession de saint Michel et du chœur admirable des Vertus célestes, que le Seigneur nous préserve du mal et nous préserve de la tentation. Amen.}
\salut{7}{Par l’intercession de saint Michel et du chœur céleste des Principautés, que Dieu remplisse nos âmes d’un véritable esprit d’obéissance. Amen.}
\salut{8}{Par l’intercession de saint Michel et du chœur céleste des Archanges, que le Seigneur nous donne la persévérance dans la foi et dans toutes les bonnes œuvres, afin que nous gagnions la gloire du paradis. Amen.}
\salut{9}{Par l’intercession de saint Michel et du chœur céleste de tous les Anges, que le Seigneur nous accorde d’être protégés par eux dans cette vie mortelle et conduits ensuite à la gloire éternelle. Amen.}

\intertitre{Antienne}
\priere Très glorieux saint Michel, chef et prince des armées célestes, gardien fidèle des âmes, vainqueur des esprits rebelles, favori de la maison de Dieu, notre admirable guide après Jésus-Christ, vous dont l’excellence et la vertu sont suréminentes : daignez nous délivrer de tous les maux, nous tous qui recourons à vous avec confiance, et faites, par votre incomparable protection, que nous avancions chaque jour dans la fidélité à servir Dieu.
\vr{Priez pour nous, ô bienheureux saint Michel, prince de l’Église de Jésus-Christ.}{Afin que nous puissions être rendus dignes de ses promesses.}

\intertitre{Oraison}
\priere Dieu tout-puissant et éternel, qui, par un prodige de bonté et de miséricorde pour le salut commun des hommes, avez choisi pour prince de votre Église le très glorieux archange saint Michel, rendez-nous dignes, nous vous en prions, d’être délivrés, par sa bienveillante protection, de tous nos ennemis, afin qu’à notre mort aucun d’eux ne puisse nous inquiéter, mais qu’il nous soit donné d’être introduits par lui en présence de votre puissante et auguste majesté. Par les mérites de Jésus-Christ Notre-Seigneur. Ainsi soit-il.

\intertitre{Prière à saint Michel Archange}
\rubrique{Du pape Léon XIII.}
\priere Saint Michel Archange, défendez-nous dans le combat ; soyez notre secours contre la malice et les embûches du démon. Que Dieu lui commande, nous vous en supplions ; et vous, Prince de la milice céleste, repoussez en enfer, par la force divine, Satan et les autres esprits mauvais qui rôdent dans le monde pour la perte des âmes. Ainsi soit-il.